\documentclass[11pt]{article}
\usepackage{amsmath,amsfonts,graphicx,geometry}
\usepackage[backend=bibtex,style=authoryear]{biblatex}
\usepackage{enumitem} 
\usepackage{graphicx}
\usepackage[utf8]{inputenc}
\usepackage[T1]{fontenc}
\usepackage{listings}
\usepackage{xcolor}
\usepackage{amssymb} 

\lstdefinestyle{mystyle}{
backgroundcolor=\color{backcolour},
commentstyle=\color{codegreen},
keywordstyle=\color{magenta},
numberstyle=\tiny\color{codegray},
stringstyle=\color{codepurple},
basicstyle=\ttfamily\footnotesize,
breakatwhitespace=false,
breaklines=true,
captionpos=b,
keepspaces=true,
numbers=left,
numbersep=5pt,
showspaces=false,
showstringspaces=false,
showtabs=false,
tabsize=2
}

\lstset{style=mystyle}

\definecolor{codegreen}{rgb}{0,0.6,0}
\definecolor{codegray}{rgb}{0.5,0.5,0.5}
\definecolor{codepurple}{rgb}{0.58,0,0.82}
\definecolor{backcolour}{rgb}{0.95,0.95,0.92}

\addbibresource{ref.bib}
\geometry{margin=1in}
\title{MTP 2: Limit order Book - Markov based approach on Market Making}
\author{Prepared by: Akhil CV, 24M1530 \\
Guided by: Prof. K.S. Mallikarjuna Rao}


\begin{document}

\maketitle

\section{Akhil's Plan }

\begin{enumerate}
\item[\checkmark] Confirm the project from Prof.
\item Identify Data foe L2/L3 levels 
\item Parse/Scape the data
\item Define the Limit orders and order flow
\item Market making problem


\subsection{Feedback from Prof: }
\item *********
\item ***********
\end{enumerate}

\maketitle

\section{Recurrence of Limit Order Book Between Price Thresholds: Theory and Empirical Verification}

\subsection{Model Setup}

Let $(X_t)_{t \in \mathbb{Z}_{\ge 0}}$ denote the discrete-time state process describing the limit order book (LOB) at time $t$.  
Let

\[
B_t = b(X_t), \qquad A_t = a(X_t),
\]

represent the best bid and best ask price functions extracted from the LOB state.

Fix thresholds
\[
L_b < U_b, \qquad L_a < U_a,
\]
and define the \emph{band region}
\[
\mathcal{B} = \left\{ x :  L_b \le b(x) \le U_b,\quad L_a \le a(x) \le U_a \right\}.
\]

Define the indicator variable of being in the band at time $t$:

\[
I_t = \mathbf{1}_{\left\{X_t \in \mathcal{B}\right\}}.
\]

Let $\tau_1$ denote the first entrance time into $\mathcal{B}$ and recursively define

\[
\tau_{k} = \inf\{t > \tau_{k-1} : X_t \in \mathcal{B}\}, \quad k \ge 2.
\]

The \emph{inter-return times} are then defined as
\[
T_k = \tau_k - \tau_{k-1}, \quad k \ge 2.
\]

\subsection{Assumption A: Stationary and Ergodic LOB Dynamics}

We work under the following standard structural assumption.

\vspace{2mm}
\noindent
\textbf{Assumption A.}
\begin{enumerate}
    \item The process $(X_t)$ is a time-homogeneous Markov chain on a measurable state space $\mathcal{X}$.

    \item The chain is irreducible and positive recurrent, and admits a unique invariant probability measure $\pi$ on $\mathcal{X}$.

    \item The chain operates in stationarity, i.e. the distribution of $X_0$ is $\pi$ (or equivalently the chain has run long enough to forget its initial state).

    \item The band region has strictly positive stationary probability:
    \[
    \pi(\mathcal{B}) = \sum_{x \in \mathcal{B}} \pi(x) > 0.
    \]
\end{enumerate}

\subsection{Definition of Recurrence}

\begin{definition}[Recurrence of the Band]
    The band $\mathcal{B}$ is called \emph{recurrent} if, starting from any state $X_0$, the LOB returns infinitely often to $\mathcal{B}$ almost surely:
    \[
    \mathbb{P}\left( X_t \in \mathcal{B} \ \text{for infinitely many } t \right) = 1.
    \]
    It is called \emph{positive recurrent} if the expected return time
    \[
    \mathbb{E}[T_2 \mid X_{\tau_1} \in \mathcal{B}] < \infty.
    \]
\end{definition}

Under Assumption A, classical Markov chain theory ensures that $\mathcal{B}$ is positive recurrent.

\subsection{Main Consistency Lemma: What the Empirical Estimators Measure}

Define the empirical time proportion spent in the band as
\[
\widehat{\pi}_n(\mathcal{B}) = \frac{1}{n} \sum_{t=1}^{n} I_t.
\]

Let $N_n = \sum_{t=1}^{n} I_t$ denote the number of visits to $\mathcal{B}$ by time $n$ and define
\[
\widehat{\mu}_m = \frac{1}{m-1}\sum_{k=2}^{m} T_k,
\]
where $m$ is the number of completed returns.

\begin{lemma}[Consistency of Empirical Recurrence Indicators]\label{lem:recurrence}
Under Assumption~A, the following statements hold.

\begin{enumerate}
    \item \textbf{Empirical time fraction consistency:}
    \[
    \widehat{\pi}_n(\mathcal{B}) \xrightarrow[n \to \infty]{a.s.} \pi(\mathcal{B}) > 0.
    \]

    \item \textbf{Consistency of mean return time:}
    The random variables $T_k$ are i.i.d.\ with finite mean and
    \[
    \widehat{\mu}_m \xrightarrow[m\to\infty]{a.s.} 
    \mathbb{E}[T_2].
    \]

    \item \textbf{Reciprocal identity:}
    In the stationary Markov chain framework
    \[
    \mathbb{E}[T_2] = \frac{1}{\pi(\mathcal{B})},
    \]
    implying that jointly
    \[
    \widehat{\mu}_m \cdot \widehat{\pi}_n(\mathcal{B}) \to 1,
    \]
    in probability under appropriate joint asymptotics.
\end{enumerate}
\end{lemma}

\subsection*{Proof Sketch}

The result follows from three standard results.

\begin{enumerate}
    \item Stationarity and ergodicity imply by the ergodic theorem:
    \[
    \frac{1}{n}\sum_{t=1}^{n}f(X_t)\to \mathbb{E}_{\pi}[f(X_0)], \quad\text{a.s.}
    \]
    Choosing $f(X_t)=\mathbf{1}_{\mathcal{B}}$ yields
    \[
    \widehat{\pi}_n(\mathcal{B}) \to \pi(\mathcal{B}).
    \]

    \item Positive recurrence ensures the existence of regeneration cycles induced by successive returns to $\mathcal{B}$, implying that the inter-return times form an i.i.d.\ sequence with finite expectation.

    \item Renewal theory for Markov chains yields
    \[
    \mathbb{E}[T_2]= \frac{1}{\pi(\mathcal{B})},
    \]
    since $\pi(\mathcal{B})$ represents the limiting frequency of visits.
\end{enumerate}

This completes the proof sketch.

\subsection{Interpretation}

The lemma shows that:

\begin{itemize}
    \item Observing a stable positive estimate of $\widehat{\pi}_n(\mathcal{B})$ provides evidence that the band has positive invariant probability.

    \item Observing a finite and stable estimate of $\widehat{\mu}_m$ provides evidence of finite return times, which characterizes positive recurrence.

    \item Agreement between $\widehat{\mu}_m$ and $1/\widehat{\pi}_n(\mathcal{B})$ supports model consistency.
\end{itemize}

Thus, empirical stability of these quantities supports the claim that the limit order book dynamics are recurrent between the selected thresholds.

\subsection{Empirical Measurement Procedure}

Given a time series of best bid and best ask data:

\[
\{(B_t, A_t) : t = 1, \dots, n\},
\]
define
\[
I_t = \mathbf{1}_{\{ L_b \le B_t \le U_b,\ L_a \le A_t \le U_a \}}.
\]

Then:

\begin{enumerate}
    \item Compute
    \[
    \widehat{\pi}_n(\mathcal{B}) = \frac{1}{n}\sum_{t=1}^{n} I_t.
    \]

    \item Identify visit times $\tau_k$ where $I_{\tau_k} = 1$, compute inter-return gaps $T_k$ and evaluate the empirical mean
    \[
    \widehat{\mu}_m = \frac{1}{m-1}\sum_{k=2}^{m} T_k.
    \]

    \item As $n$ grows, examine whether $\widehat{\pi}_n(\mathcal{B})$ converges to a positive constant and whether $\widehat{\mu}_m$ stabilizes.
\end{enumerate}

Robust convergence across trading days and instruments provides strong empirical support for recurrence.

\section{Market Regime and Jump model}

Financial markets evolve through different regimes characterized by
distinct volatility and return patterns. Examples include:

\begin{itemize}
\item Bull market regimes
\item Bear market regimes
\item High volatility periods
\item Low volatility periods
\end{itemize}

Regime-switching models aim to identify these states and detect
transitions between them.

Traditional approaches include:

\begin{itemize}
\item Hidden Markov Models
\item Gaussian Mixture Models
\item Clustering techniques
\end{itemize}

However, these models face challenges when detecting abrupt price
jumps or distinguishing them from volatility spikes.

\subsection{Jump Models}

Jump models extend clustering methods by introducing temporal
consistency through a penalty for switching regimes.

Let

\[
X = \{x_1, x_2, \dots, x_T\}
\]

be a financial return series.

We define a feature vector

\[
y_t \in \mathbb{R}^D
\]

and a hidden state sequence

\[
s_t \in \{1, \dots, K\}
\]

where $K$ represents the number of regimes.

Each state has parameters

\[
\theta_k
\]

representing the cluster centroid.

\subsection{Objective Function}

The jump model is estimated by minimizing the objective function

\[
\Psi(\theta,S) =
\sum_{t=1}^{T-1}
\left[
\ell(y_t,\theta_{s_t})
+
\lambda \eta(s_t,s_{t+1})
\right]
+
\ell(y_T,\theta_{s_T})
\]

where

\begin{itemize}
\item $\ell(y_t,\theta_{s_t})$ measures clustering error
\item $\eta(s_t,s_{t+1})$ penalizes regime switching
\item $\lambda$ controls the persistence of regimes
\end{itemize}

Typically

\[
\ell(x,y) = ||x-y||^2
\]

and

\[
\eta(s_t,s_{t+1}) = \mathbf{1}_{s_t \neq s_{t+1}}
\]

\subsection{Learning the Conditional Distribution}

Traditional jump models assume observations in a state share the
same mean. However, financial returns also vary in volatility.

To capture higher moments we assume

\[
y_t \sim \mathcal{N}(\mu_{s_t}, \sigma^2_{s_t})
\]

and estimate both mean and variance.

An extended objective function may include moment matching:

\[
(y_t - \mu_{s_t})^2 + (y_t^2 - \sigma_{s_t}^2)^2
\]

This allows the model to distinguish between volatility changes
and true price jumps.

\subsection{DeepNet Jump Models}

Clustering in the original feature space may fail when clusters
overlap. To address this, we map observations into a latent space
using a Variational Autoencoder.

Let

\[
z_t \sim q_{\phi}(z|y_t)
\]

denote the latent representation.

The VAE provides

\[
\mu_t, \sigma_t
\]

representing the approximate posterior distribution.

We construct a latent feature

\[
F_t = \mu_t + \sigma_t \epsilon_t
\]

where

\[
\epsilon_t \sim \mathcal{N}(0,I)
\]

Clustering is then performed in the latent space.

\subsection{State Sequence Estimation}

The optimal state sequence is estimated using dynamic programming.

Define

\[
L(t,k) = ||F_t - \mu_k||^2
\]

Then compute recursively

\[
V(t,k) =
L(t,k)
+
\min_j
\left[
V(t-1,j) + \lambda \mathbf{1}_{j \neq k}
\right]
\]

The optimal state sequence is obtained through backtracking.

\subsection{Detecting Price Jumps}

Price jumps are defined as sudden deviations from the underlying
distribution of returns.

Instead of detecting individual outliers, the framework evaluates
whether segments of time series are anomalous relative to their
preceding segments.

\subsection{Path Signatures}

A time series path can be represented using signatures from rough
path theory.

For a path $X_t$, the signature consists of iterated integrals

\[
S(X) =
\left(
1,
\int dX,
\int X\, dX,
\int X^2 dX,
\dots
\right)
\]

These features capture the geometry and order of the path.

\subsection{Variance Norm}

To measure deviation from a distribution we compute the variance
norm

\[
d_V(x) =
\sqrt{
(x-\mu)^T
\Sigma^{-1}
(x-\mu)
}
\]

where

\begin{itemize}
\item $\mu$ is the mean vector
\item $\Sigma$ is the covariance matrix
\end{itemize}

Large values indicate anomalous observations.

\subsection{Consecutive Signature Distance}

Given two consecutive segments $S_t$ and $S_{t+1}$ we compute

\[
d_V(S_t,S_{t+1})
\]

If the distance exceeds a threshold, a price jump is detected.

\subsection{Framework}

The full pipeline consists of:

\begin{enumerate}
\item Convert time series into overlapping segments
\item Compute path signatures
\item Calculate variance norm distances
\item Identify jumps via thresholding
\item Train jump models on filtered data
\item Predict hidden regimes
\end{enumerate}

\subsection{Simulation Results}

Experiments demonstrate that DeepNet Jump Models achieve
higher balanced accuracy compared to traditional clustering
and HMM approaches.

\begin{center}
\begin{tabular}{lcc}
\toprule
Model & Accuracy & BAC \\
\midrule
GMM & 0.379 & 0.221 \\
VAE & 0.466 & 0.367 \\
VAE + Viterbi & 0.532 & 0.454 \\
\bottomrule
\end{tabular}
\end{center}

\subsection{Conclusion}

DeepNet Jump Models provide a robust framework for regime detection
and price jump identification in financial markets.

The approach integrates:

\begin{itemize}
\item clustering-based jump models
\item deep learning representations
\item anomaly detection via variance norm
\item path signatures for time-series structure
\end{itemize}

This combination improves both predictive accuracy and responsiveness
to regime shifts.
\section{Data Processing and Market Microstructure Feature Construction}

\subsection{Market Data Representation}

Let the raw market data be represented as a sequence of limit order book observations

\[
\mathcal{D} = \{(t_i, s_i, p_i, q_i)\}_{i=1}^{N}
\]

where

\begin{itemize}
\item $t_i$ denotes the timestamp,
\item $s_i \in \{\text{Bid},\text{Ask}\}$ indicates the side of the order,
\item $p_i$ denotes the quoted price,
\item $q_i$ denotes the quoted quantity.
\end{itemize}

Each snapshot of the limit order book is indexed by a snapshot identifier

\[
\text{snap\_idx} = k
\]

which groups together all orders belonging to the same order book state.

\subsection{Midprice Construction}

For each snapshot $k$, the best bid and ask prices are defined as

\[
B_k = \max_{i:s_i=\text{Bid}} p_i
\]

\[
A_k = \min_{i:s_i=\text{Ask}} p_i
\]

The midprice is then defined as

\[
M_k = \frac{B_k + A_k}{2}.
\]

To analyze the relative structure of the book around the midprice we define a normalized price coordinate

\[
p_i^{norm} = p_i - M_k.
\]

This representation centers the order book around zero so that positive values correspond to the ask side and negative values correspond to the bid side.

\subsection{Quantity Normalization}

For derivatives instruments traded with contract size $L$, quantities are normalized as

\[
q_i^{norm} = \frac{q_i}{L}.
\]

In our experiments the NIFTY futures contract size is

\[
L = 65.
\]

This normalization allows comparability across instruments.

\section{Extraction of Level--1 Order Book}

The Level--1 order book represents the best bid and best ask at each snapshot.

Formally,

\[
B_t = \max_{i:s_i=\text{Bid}} p_i, \qquad
A_t = \min_{i:s_i=\text{Ask}} p_i.
\]

The associated quantities are

\[
Q_t^{bid}, \qquad Q_t^{ask}.
\]

The Level--1 state is therefore

\[
X_t = (B_t,A_t,Q_t^{bid},Q_t^{ask}).
\]

This representation is the primary input to the regime detection and market making models.

\section{Microstructure Feature Engineering}

From the Level--1 order book state $X_t$ we construct a set of microstructure features.

\subsection{Spread}

The bid--ask spread is

\[
S_t = A_t - B_t.
\]

\subsection{Microprice}

The microprice represents a liquidity weighted estimate of the efficient price:

\[
m_t =
\frac{A_t Q_t^{bid} + B_t Q_t^{ask}}
     {Q_t^{bid}+Q_t^{ask}}.
\]

The microprice tilts toward the side with greater queue size.

\subsection{Microprice Return}

The return process is defined as

\[
r_t = \log\left(\frac{m_t}{m_{t-1}}\right).
\]

\subsection{Order Book Imbalance}

Order flow imbalance is defined as

\[
I_t =
\frac{Q_t^{bid} - Q_t^{ask}}
     {Q_t^{bid}+Q_t^{ask}}.
\]

\subsection{Volume Pressure}

To capture asymmetry in liquidity supply we define

\[
V_t = \log\left(\frac{Q_t^{bid}+1}{Q_t^{ask}+1}\right).
\]

The complete feature vector is therefore

\[
Y_t =
(r_t,S_t,I_t,V_t).
\]

\section{Hidden Markov Model for Market Regime Detection}

Financial markets operate under multiple regimes such as market making periods and directional order flow phases.

We model the regime sequence using a Hidden Markov Model (HMM).

\subsection{Latent State Process}

Let

\[
S_t \in \{1,2,3\}
\]

represent the hidden regime at time $t$.

We assume the regime evolves as a Markov chain

\[
P(S_t=j|S_{t-1}=i) = P_{ij}.
\]

\subsection{Observation Model}

Conditional on the state $S_t=k$, the feature vector follows a Gaussian distribution

\[
Y_t | S_t=k \sim
\mathcal{N}(\mu_k,\Sigma_k).
\]

The parameters

\[
(\mu_k,\Sigma_k)
\]

are estimated via maximum likelihood using the Baum--Welch algorithm.

\subsection{Interpretation of Regimes}

After estimation the states are interpreted as

\begin{itemize}
\item $S_t=0$: Market making regime
\item $S_t=1$: Upward pressure regime
\item $S_t=2$: Downward pressure regime
\end{itemize}

These regimes correspond to different liquidity conditions in the order book.

\section{Markov Limit Order Book and Mean Reversion}

During market making periods the microprice exhibits mean reverting behavior.

We model the microprice using an Ornstein--Uhlenbeck process

\[
dm_t = -\kappa(m_t-\mu)dt + \sigma dW_t.
\]

Here

\begin{itemize}
\item $\mu$ is the equilibrium price,
\item $\kappa$ is the mean reversion speed,
\item $\sigma$ is the volatility parameter.
\end{itemize}

\subsection{Discrete Representation}

For discrete observations

\[
m_{t+1} = a m_t + b + \epsilon_t
\]

with

\[
\epsilon_t \sim \mathcal{N}(0,\sigma^2).
\]

The parameters are estimated using linear regression.

\subsection{Stationary Distribution}

The stationary distribution of the OU process is

\[
m_t \sim
\mathcal{N}\left(
\mu,
\frac{\sigma^2}{2\kappa}
\right).
\]

The equilibrium mean is

\[
\mu = \frac{b}{1-a}.
\]

The stationary variance is

\[
\sigma_{stat}^2 = \frac{\sigma^2}{1-a^2}.
\]

\section{Market Making Price Bands}

Using the stationary distribution we define upper and lower microprice bounds

\[
L_t = \mu - k\sigma_{stat},
\]

\[
U_t = \mu + k\sigma_{stat}.
\]

The constant $k$ determines the width of the band.

\section{Conversion to Bid and Ask Bounds}

Since quoting occurs at bid and ask prices rather than microprice, the band is translated using the spread.

Let

\[
S_t = A_t - B_t.
\]

The quoting bounds become

\[
B_t^{bound} = L_t - \frac{S_t}{2}
\]

\[
A_t^{bound} = U_t + \frac{S_t}{2}.
\]

\section{Market Microstructure Constraints}

To maintain a valid order book we enforce

\[
B_t^{bound} \le B_t \le A_t \le A_t^{bound}.
\]

Violations are corrected by clipping

\[
B_t^{bound} = \min(B_t^{bound},B_t)
\]

\[
A_t^{bound} = \max(A_t^{bound},A_t).
\]

\section{Algorithmic Pipeline}

The full estimation procedure consists of the following steps.

\begin{enumerate}
\item Read raw order book snapshots.
\item Compute midprice and normalize prices.
\item Extract Level--1 order book states.
\item Construct microstructure features.
\item Train the Hidden Markov Model.
\item Infer market regimes.
\item During market making regimes estimate OU parameters.
\item Compute stationary microprice bands.
\item Convert microprice bands to bid and ask bounds.
\end{enumerate}

\section{Interpretation}

The resulting bands represent the equilibrium region in which market makers are willing to provide liquidity.

When the microprice approaches the upper bound, the order book tends to experience selling pressure.

When the microprice approaches the lower bound, buying pressure increases.

Thus the model provides a quantitative representation of market making behavior in the limit order book.


% \subsection{Python Implementation (Appendix)}

% \begin{verbatim}
% import numpy as np
% import pandas as pd

% def compute_band_recurrence_stats(df, L_b, U_b, L_a, U_a):
%     in_band = (
%         (df['bid'] >= L_b) & (df['bid'] <= U_b) &
%         (df['ask'] >= L_a) & (df['ask'] <= U_a)
%     )

%     visit_times = df.index[in_band]

%     if len(visit_times) < 2:
%         inter_return_steps = np.array([])
%         inter_return_seconds = np.array([])
%     else:
%         visit_idx = np.where(in_band.values)[0]
%         inter_return_steps = np.diff(visit_idx)

%         inter_return_seconds = np.diff(
%             visit_times.values.astype('datetime64[ns]')
%         ) / np.timedelta64(1, 's')

%     time_fraction = in_band.mean()

%     time_fraction_over_time = in_band.cumsum() / np.arange(1, len(in_band)+1)
%     time_fraction_over_time.index = df.index

%     return {
%         'in_band': in_band,
%         'visit_times': visit_times,
%         'inter_return_steps': inter_return_steps,
%         'inter_return_seconds': inter_return_seconds,
%         'time_fraction': time_fraction,
%         'time_fraction_over_time': time_fraction_over_time
%     }
% \end{verbatim}




\printbibliography  % This will print the bibliography at the end of your document
\end{document}
